\section*{Bab \ref{ch:g5:division} --- Pembagian dan Kelipatan}

\begin{solution}{exo:g5:division:1}
$851 \div 23$: \omterm{def:g5:division:multiple}{kelipatan} $23$: $23, 46, 69, 92, 115, 138, 161,
184, 207$. Lalu $85 \div 23$: $3$ kali ($69$), \omterm{def:g3:division:remainder}{sisa} $16$;
turunkan angka $1$: $161 \div 23 = 7$ tepat. \omterm{def:g3:division:remainder}{Hasil bagi} $37$,
\omterm{def:g3:division:remainder}{sisa} $0$ ($23 \times 37 = 851$).

$704 \div 32$: $70 \div 32$: $2$ kali ($64$), \omterm{def:g3:division:remainder}{sisa} $6$; turunkan
angka $4$: $64 \div 32 = 2$. \omterm{def:g3:division:remainder}{Hasil bagi} $22$, \omterm{def:g3:division:remainder}{sisa} $0$.
\end{solution}

\begin{solution}{exo:g5:division:2}
$9 \div 2 = 4.5$; \quad $27 \div 4 = 6.75$; \quad
$33 \div 5 = 6.6$; \quad $21 \div 8 = 2.625$.
\end{solution}

\begin{solution}{exo:g5:division:3}
$54 \div 4 = 13.5$: masing-masing membayar $13.50$.
\end{solution}

\begin{solution}{exo:g5:division:4}
$10 \div 3 = 3.333\dots$: angka $3$ berulang selamanya ---
pembagiannya tidak pernah berhenti. \omterm{def:g4:numbers:round}{Dibulatkan} ke perseratusan: $3.33$.
\end{solution}

\begin{solution}{exo:g5:division:5}
\omterm{def:g5:division:multiple}{Kelipatan} $7$ sampai $70$: $7, 14, 21, 28, 35, 42, 49, 56, 63,
70$. \omterm{def:g5:division:multiple}{Pembagi} dari $18$: $1, 2, 3, 6, 9, 18$. \omterm{def:g5:division:multiple}{Pembagi} dari $24$:
$1, 2, 3, 4, 6, 8, 12, 24$.
\end{solution}

\begin{solution}{exo:g5:division:6}
$470$: habis dibagi $2$ (berakhir $0$), oleh $5$, oleh $10$; \omterm{def:g1:addition:def}{jumlah}
angkanya $11$: tidak habis dibagi $3$.

$735$: berakhir $5$: habis dibagi $5$, tidak oleh $2$ maupun $10$;
\omterm{def:g1:addition:def}{jumlah} angkanya $15$: habis dibagi $3$.

$8\,001$: \omterm{def:g3:numbers:evenodd}{ganjil}; \omterm{def:g1:addition:def}{jumlah} angkanya $9$: hanya habis dibagi $3$.

$1\,314$: \omterm{def:g3:numbers:evenodd}{genap}; \omterm{def:g1:addition:def}{jumlah} angkanya $9$: habis dibagi $2$ dan $3$
(karena itu juga oleh $6$), tidak oleh $5$.
\end{solution}

\begin{solution}{exo:g5:division:7}
Habis dibagi $3$ dan $5$ berarti habis dibagi $15$: antara $60$ dan
$80$, yaitu bilangan $60$ dan $75$.
\end{solution}

\begin{solution}{exo:g5:division:8}
$7.2 \div 6 = 1.2$: setiap bagian panjangnya $1.2$ m (periksa:
$1.2 \times 6 = 7.2$).
\end{solution}

\begin{solution}{exo:g5:division:9}
$1\,000 = 24 \times 41 + 16$: empat puluh satu kantong penuh, dan
$16$ kelereng tersisa.
\end{solution}

\begin{solution}{exo:g5:division:10}
$7 \div 5 = 1.4$: masing-masing dari $5$ teman itu mendapat $1.4$
batang --- mungkin dilakukan dengan memotong batang menjadi
persepuluhan. Untuk $5 \div 7$: pembagiannya memberi $0.714285\dots$,
yang tidak pernah berhenti --- tidak ada bagian desimal yang tepat;
hanya \omterm{def:g4:fractions:def}{pecahan} $\frac57$ yang menamainya dengan tepat.
\end{solution}

\begin{solution}{exo:g5:division:11}
\omterm{def:g1:addition:def}{Jumlah} angkanya: $5 + 2 + d = 7 + d$. Habis dibagi $3$ ketika
$7 + d$ sama dengan $9$, $12$ atau $15$: $d = 2$, $5$ atau $8$. Yang
terkecil adalah $d = 2$ (memberi $522 = 3 \times 174$).

\end{solution}
